%!TeX root=../pridetop.tex
\chapter[Chapter \thechapter]{}
	
\begin{figure}[t!]
\centering
\includegraphics[width=.95\linewidth]{29top}
\captionlistentry{<Lady Catherine,> said she, <you have given me a treasure.>}
\end{figure}


\lettrine[lines=6,image=true]{initials/chap29m}{r}  Collins's triumph, in consequence of this invitation, was complete. The power of displaying the grandeur of his patroness to his wondering visitors, and of letting them see her civility towards himself and his wife, was exactly what he had wished for; and that an opportunity of doing it should be given so soon was such an instance of Lady Catherine's condescension as he knew not how to admire enough.

<I confess,> said he, <that I should not have been at all surprised by her Ladyship's asking us on Sunday to drink tea and spend the evening at Rosings. I rather expected, from my knowledge of her affability, that it would happen. But who could have foreseen such an attention as this? Who could have imagined that we should receive an invitation to dine there (an invitation, moreover, including the whole party) so immediately after your arrival?>

<I am the less surprised at what has happened,> replied Sir William, <from that knowledge of what the manners of the great really are, which my situation in life has allowed me to acquire. About the court, such instances of elegant breeding are not uncommon.>

Scarcely anything was talked of the whole day or next morning but their visit to Rosings. Mr Collins was carefully instructing them in what they were to expect, that the sight of such rooms, so many servants, and so splendid a dinner, might not wholly overpower them.

When the ladies were separating for the toilette, he said to Eliz\-a\-beth,—

<Do not make yourself uneasy, my dear cousin, about your apparel. Lady Catherine is far from requiring that elegance of dress in us which becomes herself and daughter. I would advise you merely to put on whatever of your clothes is superior to the rest—there is no occasion for anything more. Lady Catherine will not think the worse of you for being simply dressed. She likes to have the distinction of rank preserved.>

While they were dressing, he came two or three times to their different doors, to recommend their being quick, as Lady Catherine very much objected to be kept waiting for her dinner. Such formidable accounts of her Ladyship, and her manner of living, quite frightened Maria Lucas, who had been little used to company; and she looked forward to her introduction at Rosings with as much apprehension as her father had done to his presentation at St James's.

As the weather was fine, they had a pleasant walk of about half a mile across the park. Every park has its beauty and its prospects; and Elizabeth saw much to be pleased with, though she could not be in such raptures as Mr Collins expected the scene to inspire, and was but slightly affected by his enumeration of the windows in front of the house, and his relation of what the glazing altogether had originally cost Sir Lewis de Bourgh.

When they ascended the steps to the hall, Maria's alarm was every moment increasing, and even Sir William did not look perfectly calm. Elizabeth's courage did not fail her. She had heard nothing of Lady Catherine that spoke her awful from any extraordinary talents or miraculous virtue, and the mere stateliness of money and rank she thought she could witness without trepidation.

From the entrance hall, of which Mr Collins pointed out, with a rapturous air, the fine proportion and finished ornaments, they followed the servants through an antechamber to the room where Lady Catherine, her daughter, and Mrs Jenkinson were sitting. Her Ladyship, with great condescension, arose to receive them; and as Mrs Collins had settled it with her husband that the office of introduction should be hers, it was performed in a proper manner, without any of those apologies and thanks which he would have thought necessary.

In spite of having been at St James's, Sir William was so completely awed by the grandeur surrounding him, that he had but just courage enough to make a very low bow, and take his seat without saying a word; and his daughter, frightened almost out of her senses, sat on the edge of her chair, not knowing which way to look. Elizabeth found herself quite equal to the scene, and could observe the three ladies before her composedly. Lady Catherine was a tall, large woman, with strongly-marked features, which might once have been handsome. Her air was not conciliating, nor was her manner of receiving them such as to make her visitors forget their inferior rank. She was not rendered formidable by silence: but whatever she said was spoken in so authoritative a tone as marked her self-importance, and brought Mr Wickham immediately to Elizabeth's mind; and, from the observation of the day altogether, she believed Lady Catherine to be exactly what he had represented.

When, after examining the mother, in whose countenance and deportment she soon found some resemblance of Mr Darcy, she turned her eyes on the daughter, she could almost have joined in Maria's astonishment at her being so thin and so small. There was neither in figure nor face any likeness between the ladies. Miss de Bourgh was pale and sickly: her features, though not plain, were insignificant; and she spoke very little, except in a low voice, to Mrs Jenkinson, in whose appearance there was nothing remarkable, and who was entirely engaged in listening to what she said, and placing a screen in the proper direction before her eyes.

After sitting a few minutes, they were all sent to one of the windows to admire the view, Mr Collins attending them to point out its beauties, and Lady Catherine kindly informing them that it was much better worth looking at in the summer.

The dinner was exceedingly handsome, and there were all the servants, and all the articles of plate which Mr Collins had promised; and, as he had likewise foretold, he took his seat at the bottom of the table, by her Ladyship's desire, and looked as if he felt that life could furnish nothing greater. He carved and ate and praised with delighted alacrity; and every dish was commended first by him, and then by Sir William, who was now enough recovered to echo whatever his son-in-law said, in a manner which Elizabeth wondered Lady Catherine could bear. But Lady Catherine seemed gratified by their excessive admiration, and gave most gracious smiles, especially when any dish on the table proved a novelty to them. The party did not supply much conversation. Elizabeth was ready to speak whenever there was an opening, but she was seated between Charlotte and Miss de Bourgh—the former of whom was engaged in listening to Lady Catherine, and the latter said not a word to her all the dinnertime. Mrs Jenkinson was chiefly employed in watching how little Miss de Bourgh ate, pressing her to try some other dish and fearing she was indisposed. Maria thought speaking out of the question, and the gentlemen did nothing but eat and admire.

When the ladies returned to the drawing-room, there was little to be done but to hear Lady Catherine talk, which she did without any intermission till coffee came in, delivering her opinion on every subject in so decisive a manner as proved that she was not used to have her judgment controverted. She inquired into Charlotte's domestic concerns familiarly and minutely, and gave her a great deal of advice as to the management of them all; told her how everything ought to be regulated in so small a family as hers, and instructed her as to the care of her cows and her poultry. Elizabeth found that nothing was beneath this great lady's attention which could furnish her with an occasion for dictating to others. In the intervals of her discourse with Mrs Collins, she addressed a variety of questions to Maria and Elizabeth, but especially to the latter, of whose connections she knew the least, and who, she observed to Mrs Collins, was a very genteel, pretty kind of girl. She asked her at different times how many sisters she had, whether they were older or younger than herself, whether any of them were likely to be married, whether they were handsome, where they had been educated, what carriage her father kept, and what had been her mother's maiden name? Elizabeth felt all the impertinence of her questions, but answered them very composedly. Lady Catherine then observed,—

<Your father's estate is entailed on Mr Collins, I think? For your sake,> turning to Charlotte, <I am glad of it; but otherwise I see no occasion for entailing estates from the female line. It was not thought necessary in Sir Lewis de Bourgh's family. Do you play and sing, Miss Bennet?>

<A little.>

<Oh then—some time or other we shall be happy to hear you. Our instrument is a capital one, probably superior to — you shall try it some day. Do your sisters play and sing?>

<One of them does.>

<Why did not you all learn? You ought all to have learned. The Miss Webbs all play, and their father has not so good an income as yours. Do you draw?>

<No, not at all.>

<What, none of you?>

<Not one.>

<That is very strange. But I suppose you had no opportunity. Your mother should have taken you to town every spring for the benefit of masters.>

<My mother would have no objection, but my father hates London.>

<Has your governess left you?>

<We never had any governess.>

<No governess! How was that possible? Five daughters brought up at home without a governess! I never heard of such a thing. Your mother must have been quite a slave to your education.>

Elizabeth could hardly help smiling, as she assured her that had not been the case.

<Then who taught you? who attended to you? Without a governess, you must have been neglected.>

<Compared with some families, I believe we were; but such of us as wished to learn never wanted the means. We were always encouraged to read, and had all the masters that were necessary. Those who chose to be idle certainly might.>

<Why, no doubt: but that is what a governess will prevent; and if I had known your mother, I should have advised her most strenuously to engage one. I always say that nothing is to be done in education without steady and regular instruction, and nobody but a governess can give it. It is wonderful how many families I have been the means of supplying in that way. I am always glad to get a young person well placed out. Four nieces of Mrs Jenkinson are most delightfully situated through my means; and it was but the other day that I recommended another young person, who was merely accidentally mentioned to me, and the family are quite delighted with her. Mrs Collins, did I tell you of Lady Metcalfe's calling yesterday to thank me? She finds Miss Pope a treasure. <Lady Catherine,> said she, <you have given me a treasure.> Are any of your younger sisters out, Miss Bennet?>

<Yes, ma'am, all.>

<All! What, all five out at once? Very odd! And you only the second. The younger ones out before the elder are married! Your younger sisters must be very young?>

<Yes, my youngest is not sixteen. Perhaps \textit{she} is full young to be much in company. But really, ma'am, I think it would be very hard upon younger sisters that they should not have their share of society and amusement, because the elder may not have the means or inclination to marry early. The last born has as good a right to the pleasures of youth as the first. And to be kept back on \textit{such} a motive! I think it would not be very likely to promote sisterly affection or delicacy of mind.>

<Upon my word,> said her Ladyship, <you give your opinion very decidedly for so young a person. Pray, what is your age?>

<With three younger sisters grown up,> replied Elizabeth, smiling, <your Ladyship can hardly expect me to own it.>

Lady Catherine seemed quite astonished at not receiving a direct answer; and Elizabeth suspected herself to be the first creature who had ever dared to trifle with so much dignified impertinence.

<You cannot be more than twenty, I am sure,—therefore you need not conceal your age.>

<I am not one-and-twenty.>

When the gentlemen had joined them, and tea was over, the card tables were placed. Lady Catherine, Sir William, and Mr and Mrs Collins sat down to quadrille; and as Miss de Bourgh chose to play at cassino, the two girls had the honour of assisting Mrs Jenkinson to make up her party. Their table was superlatively stupid. Scarcely a syllable was uttered that did not relate to the game, except when Mrs Jenkinson expressed her fears of Miss de Bourgh's being too hot or too cold, or having too much or too little light. A great deal more passed at the other table. Lady Catherine was generally speaking—stating the mistakes of the three others, or relating some anecdote of herself. Mr Collins was employed in agreeing to everything her Ladyship said, thanking her for every fish he won, and apologizing if he thought he won too many. Sir William did not say much. He was storing his memory with anecdotes and noble names.

When Lady Catherine and her daughter had played as long as they chose, the tables were broken up, the carriage was offered to Mrs Collins, gratefully accepted, and immediately ordered. The party then gathered round the fire to hear Lady Catherine determine what weather they were to have on the morrow. From these instructions they were summoned by the arrival of the coach; and with many speeches of thankfulness on Mr Collins's side, and as many bows on Sir William's, they departed. As soon as they had driven from the door, Elizabeth was called on by her cousin to give her opinion of all that she had seen at Rosings, which, for Charlotte's sake, she made more favourable than it really was. But her commendation, though costing her some trouble, could by no means satisfy Mr Collins, and he was very soon obliged to take her Ladyship's praise into his own hands.
